\section{对标中国船舶：差异不在“有没有预告”，而在能否建立经营到 EPS 的桥}

中国船舶案例提供的正确参照不是把报告页数、目标价或评级机械搬到雅化，而是先问：在正式半年报前，已有公开资料能否把\textbf{可观察经营变量}一路连接到收入、利润、现金和每股收益。中国船舶同样处在 H1 预告而非正式中报阶段，但其上市公司口径已经披露了重组后股本、合并范围、在手合同、年度交付计划、合同负债和历史成本/现金流；这些资料允许其以合并口径做受限的交付--毛利--现金情景，并将未披露的逐船 ASP、军品和体外资产明确设为零信用。

雅化目前的结论不同，并\textbf{不是}因为“没有 H1 正式报告”这一项本身。雅化有已审计的 2025 年锂盐销量、分部收入/毛利和民爆历史底座，也有 H1 利润预告；但没有 2026H1 的锂盐销量、实际 ASP、原料/加工成本、自供比例、分部利润和经营现金流。锂盐是短周期、销售确认型业务，缺少这些字段时，行业现货价、设计产能和客户资格不能代替收入--毛利--现金桥。把中国船舶的目标价形式复制到雅化，只会重新制造“有估值表、没有估值”的问题。

\begin{exhibitbox}[完整覆盖的门槛：同为 H1 前，证据结构不同]
\begin{tabularx}{\exhibitboxwidth}{L{2.6cm}X X X}
\toprule
门槛 & 中国船舶案例可用的输入 & 雅化截至 7 月 23 日的输入 & 对雅化模型的结论 \\
\midrule
经营可见性 & 上市体在手订单、交付计划、重组后范围和合同负债可构成收入约束 & 只有 2025 年销量/分部历史与 H1 利润区间；没有 H1 实际销量、售价或订单覆盖 & 不把设计产能、客户名称或行业现货价转成 H2 收入 \\
利润引擎 & 交付、产品结构、毛利率、确认率和营运资金可做合并情景敏感性 & 锂盐量、实现价、原料/加工成本及民爆 2026 利润未披露 & 不建立自有 2026E/2027E EPS，也不使用单季年化 \\
现金验证 & 合同负债、合同资产、存货和年度现金流能约束订单兑现质量 & 2025A、2026Q1 已显示现金转换压力；H1 现金流尚未披露 & 低前瞻 PE 不能被解释为“便宜” \\
股本与估值 & 合并后股本、经营模型与周期正常化方法能够相互校验 & 股本和市值可复算，但公平价值分母缺失 & 只保留现价对外部 EPS 的交叉核验 \\
未披露项处理 & 军品、体外订单、逐船 ASP 均为零估值信用或显式敏感性 & 资源成本优势、客户订单、硫化锂、民爆 H1 增量同样均为零信用 & 不用“潜在弹性”填补模型缺口 \\
\bottomrule
\end{tabularx}
\end{exhibitbox}
\sourcenote{方法参照：中国船舶（600150.SH）完整研究案例的 Research Brief、Artifact Contract、Growth Earnings Model、Valuation Model 与 Final Sign-off（数据截至 2026-07-22）；雅化集团 2025 年年报、2026Q1 报告及 2026H1 业绩预告。中国船舶案例仅用于研究方法比较，不向雅化转移任何公司事实或估值参数。}

\section{H2 分歧被量化了：预告证明 Q2，不替代下半年的持续性}

已披露的 H1 归母 11.00--13.00 亿元，覆盖三家外部全年归母预测的 36.3\%--56.2\%。这说明 H1 的利润事件很强，但仍有约一半全年利润需要 H2 实现。更重要的是，三家报告要求的 H2 单季平均利润差异很大：国信的要求明显低于预告 Q2，开源仍需维持较高水位，东吴则要求下半年两季基本维持或超过预告 Q2 的上沿。分歧不在 H1 是否改善，而在 H2 的量、价、成本和现金能否共同延续。

\begin{exhibitbox}[外部 H2 要求相对预告 Q2：不是概率预测，而是验证强度尺]
\begin{tabularx}{\exhibitboxwidth}{L{2.4cm}R{2.7cm}R{2.7cm}X}
\toprule
机构 & 外部 H2 单季平均归母 & 相对 Q2 预告区间的端点比例 & 对正式 H1 的验证要求 \\
\midrule
国信证券 & 5.08--6.08 亿元 & 约为 Q2 低/高端的 66.7\% / 63.3\% & H1 需证明 Q2 并非依赖不可持续一次性项目；H2 可允许季节性回落 \\
开源证券 & 6.35--7.35 亿元 & 约为 Q2 低/高端的 83.4\% / 76.5\% & 需看到销量、毛利和营运资本均改善，才能把 Q2 视为更高盈利中枢 \\
东吴证券 & 8.66--9.66 亿元 & 约为 Q2 低/高端的 113.8\% / 100.5\% & 需要连续两个季度接近 Q2 高位；没有公司级量价成本和现金桥时，不给予预先信用 \\
\bottomrule
\end{tabularx}
\end{exhibitbox}
\sourcenote{国信、开源、东吴证券原始报告（2026-04-29 至 2026-07-07）；雅化集团《2026 年半年度业绩预告》（2026-07-07）。H2 单季平均和比例均为本研究以各机构全年归母预测、H1 预告及公司披露 Q2 区间作出的算术推导；不代表雅化指引或本报告预测。}

这给出当前最有信息含量的判断：\textbf{市场并不需要再确认“Q2 有没有弹性”，而需要确认“高利润是否不用继续扩大应收、存货和原料占用，就能跨越 H2”。} 在正式 H1 之前，把东吴式高 H2 假设当作已发生，或把国信式温和 H2 假设当作安全垫，都是越过证据边界；二者都只能作为下一次模型的交叉校验。

\section{可反驳的核心判断与最强反方}

本事件研究不是“没有观点”。核心判断是：\textbf{H1 改善已被预告确认，但现有披露更充分支持行业价格/需求 $\beta$，不足以量化雅化自身的成本、资源和民爆缓冲对 H2 的独立贡献。} 因而，利润质量和 H2 持续性仍是现价含义的决定变量。该判断可被正式 H1 直接推翻：若收入、锂盐量价成本、民爆分部和经营现金流同时显示改善，则应从事件件升级为完整盈利模型；若利润落在预告区间但现金/营运资本恶化，当前的低 PE 交叉值反而更接近周期折价而非价值缺口。

最强反方是，雅化的矿--产--销协同、资源安排、头部客户与民爆高毛利底座已经使本轮盈利具有较强公司特异性。这个反方不是被否定，而是暂未被量化：它需要实际供矿/自给、单位成本、产品结构、分部利润和回款来取得估值信用。研究纪律要求把未验证的优势定义为\textbf{上行验证项}，而非先塞进基准 EPS。
